\chapter{Programs — List}
%%% generated by the blueprint pipeline from TCSlib.Programs.Util.List
\section{Overview}

Declarations of the Lean module \texttt{TCSlib.Programs.Util.List}.

\section{Declarations}

\begin{theorem}[Predicate preserved by defaulted list access]
\lean{List.getD_of_forall_mem}
\label{List.getD_of_forall_mem}
\leanok
\difficulty{4}
Let $p$ be a predicate on a type, let $l$ be a finite list of elements of that type, and
let $d$ be a default value. If $p$ holds of $d$ and of every entry of $l$, then $p$
holds of the entry selected from $l$ at any natural-number index, where out-of-range
access returns $d$.
\end{theorem}

\begin{theorem}[Lookup in an appended list]
\lean{List.getD_append_add}
\label{List.getD_append_add}
\leanok
\difficulty{3}
Let $l$ and $l'$ be lists with entries in a type $\alpha$, let $j$ be a natural number,
and let $d\in\alpha$ be a default value. The lookup at position
$\operatorname{length}(l)+j$ in the concatenation $l++l'$, using default $d$, equals the
lookup at position $j$ in $l'$, using the same default.
\end{theorem}

\begin{theorem}[Default lookup after dropping a prefix]
\lean{List.getD_drop}
\label{List.getD_drop}
\leanok
\difficulty{1}
For a list $l$, natural numbers $a,j$, and default element $d$, the lookup at position
$j$ in the list obtained by dropping the first $a$ entries of $l$, with default $d$,
equals the lookup at position $a+j$ in $l$ with the same default.
\end{theorem}

\begin{theorem}[Access within a truncated list]
\lean{List.getD_take_of_lt}
\label{List.getD_take_of_lt}
\leanok
\difficulty{1}
Let $l$ be a finite list, let $n$ be a natural number, and let $j<n$. For any default
value $d$, the entry at index $j$ in the list formed from the first $n$ entries of $l$
equals the entry at index $j$ in $l$.
\end{theorem}

\begin{theorem}[Indexing a mapped range]
\lean{List.getD_map_range}
\label{List.getD_map_range}
\leanok
\difficulty{3}
For a function $f:\mathbb{N}\to\alpha$, a natural number $n$, an index $i<n$, and a
default value $d\in\alpha$, the $i$th entry of the list obtained by mapping $f$ over
$0,1,\ldots,n-1$ is $f(i)$.
\end{theorem}

\begin{theorem}[Lookup in a mapped list]
\lean{List.getD_map_of_lt}
\label{List.getD_map_of_lt}
\leanok
\difficulty{3}
Let $f$ map elements of a type $\alpha$ to elements of a type $\beta$, and let $l$ be a
finite list of elements of $\alpha$. For any position $i<|l|$ and defaults $d\in\alpha$
and $d'\in\beta$, the lookup at position $i$ in the list obtained by applying $f$ to
every entry of $l$, with default $d'$, equals $f$ applied to the lookup at position $i$
in $l$ with default $d$.
\end{theorem}

\begin{theorem}[Access after changing a list entry]
\lean{List.getD_set_of_lt}
\label{List.getD_set_of_lt}
\leanok
\difficulty{4}
Let $l$ be a finite list, let $q$ be an index with $q<\operatorname{length}(l)$, and let
$a,d$ be elements of the same type. At any index $q'$, the entry at index $q'$ in the
list obtained by replacing the entry at $q$ with $a$ is $a$ if $q'=q$, and otherwise is
the entry at $q'$ in $l$; out-of-range access uses the default value $d$.
\end{theorem}

\begin{theorem}[Default lookup after appending one element]
\lean{List.getD_append_singleton}
\label{List.getD_append_singleton}
\leanok
\difficulty{4}
For a list $l$ and elements $a,d$ of its type, the function that looks up position $e$
in the list formed by appending $a$ to $l$, using default $d$, is the function that
looks up position $e$ in $l$ with the same default, except that its value at position
$|l|$ is replaced by $a$.
\end{theorem}

\begin{theorem}[Initial segment ending at a specified entry]
\lean{List.take_succ_getD}
\label{List.take_succ_getD}
\leanok
\difficulty{3}
Let $l$ be a list, let $i$ be a natural number less than the length of $l$, and let $d$
be an element of the same type as the entries of $l$. The first $i+1$ entries of $l$ are
the first $i$ entries followed by the entry at position $i$.
\end{theorem}

\begin{theorem}[Lookup in a concatenation with varying lengths]
\lean{List.getD_flatMap_range_sum}
\label{List.getD_flatMap_range_sum}
\leanok
\uses{List.getD_append_add}
\difficulty{4}
Let $f$ assign a list to each natural number, and let $n,r,o$ be natural numbers with
$r<n$ and $o<|f(r)|$. For a default element $d$, the entry at position equal to the sum
of the lengths of $f(0),\ldots,f(r-1)$ plus $o$ in the concatenation of
$f(0),\ldots,f(n-1)$, with default $d$, equals the entry at position $o$ in $f(r)$ with
the same default.
\end{theorem}

\begin{theorem}[Length of a flattened range]
\lean{List.length_flatMap_range}
\label{List.length_flatMap_range}
\leanok
\difficulty{3}
Let $n$ and $k$ be natural numbers, and let $f$ assign a finite list to each natural
number. If $\text{length}(f(z))=k$ for every natural number $z<n$, then concatenating
the lists $f(0),\ldots,f(n-1)$ gives a list of length $nk$.
\end{theorem}

\begin{theorem}[Lookup in a concatenation over a range]
\lean{List.getD_flatMap_range}
\label{List.getD_flatMap_range}
\leanok
\uses{List.getD_flatMap_range_sum, List.length_flatMap_range}
\difficulty{3}
Let $k,n$ be natural numbers and let $f$ assign a list to each natural number, with
$f(z)$ of length $k$ whenever $z<n$. For $z<n$, $r<k$, and a default element $d$, the
entry at position $zk+r$ in the concatenation of $f(0),\ldots,f(n-1)$, with default $d$,
equals the entry at position $r$ in $f(z)$, with the same default.
\end{theorem}

\begin{theorem}[Length of a flat-mapped finite range]
\lean{List.length_flatMap_finRange}
\label{List.length_flatMap_finRange}
\leanok
\difficulty{1}
Let $n,k$ be natural numbers and assign to each index $i<n$ a finite list, each of
length $k$. The list formed by concatenating these lists in index order has length $n
k$.
\end{theorem}

\begin{theorem}[Lookup in a concatenation of equal-length lists]
\lean{List.getD_flatMap_finRange}
\label{List.getD_flatMap_finRange}
\leanok
\uses{List.getD_flatMap_range, List.map_coe_finRange_eq_range}
\difficulty{4}
Let $n,k$ be natural numbers and let $f$ assign a list to each natural number less than
$n$, with every such list having length $k$. For $i<n$, a natural number $r<k$, and a
default element $d$, the entry at position $ik+r$ in the concatenation of these lists,
with default $d$, equals the entry at position $r$ in the list assigned to $i$, with the
same default.
\end{theorem}

\begin{theorem}[Extracting a block from concatenated lists]
\lean{List.take_drop_flatMap_range}
\label{List.take_drop_flatMap_range}
\leanok
\uses{List.length_flatMap_range}
\difficulty{4}
Let $k,m,n,c$ be natural numbers and let $f$ assign a list to each natural number.
Suppose every list $f(z)$ has length $k$ whenever $z<m+n+c$. From the concatenation of
$f(0),\ldots,f(m+n+c-1)$, discard the first $mk$ entries and retain the next $nk$
entries. The resulting list is the concatenation of $f(m),\ldots,f(m+n-1)$.
\end{theorem}

\begin{theorem}[Pairwise mapping of concatenated lists]
\lean{List.zipWith_flatMap_range}
\label{List.zipWith_flatMap_range}
\leanok
\uses{List.length_flatMap_range}
\difficulty{4}
Let $g$ be a function of two arguments, let $n,k$ be natural numbers, and let $f$ and
$f'$ assign lists to each natural number, with entries in the respective input types of
$g$. Suppose every list $f(z)$ and every list $f'(z)$ has length $k$. Applying $g$ entry
by entry to the concatenations of $f(0),\ldots,f(n-1)$ and $f'(0),\ldots,f'(n-1)$ gives
the concatenation, over $z<n$, of the lists obtained by applying $g$ entry by entry to
$f(z)$ and $f'(z)$.
\end{theorem}

\begin{theorem}[Length bound for a flattened list]
\lean{List.length_flatMap_le}
\label{List.length_flatMap_le}
\leanok
\difficulty{3}
Let $l$ be a finite list and let $f$ assign a finite list to each element of its type.
If every element $a$ of $l$ satisfies $\text{length}(f(a))\le c$ for some natural number
$c$, then the length of the list obtained by concatenating the lists $f(a)$ over $a\in
l$ is at most $\text{length}(l)c$.
\end{theorem}

\begin{theorem}[Sum over an initial range]
\lean{List.sum_map_range}
\label{List.sum_map_range}
\leanok
\difficulty{1}
Let $M$ be a commutative additive monoid, let $f:\bbn\to M$ be a function on the natural
numbers, and let $n\in\bbn$. The sum of the list obtained by applying $f$ to $0,1,\ldots,n-1$
equals the sum of $f(i)$ over all natural numbers $i<n$.
\end{theorem}

\begin{theorem}[Monotonicity of initial range sums]
\lean{List.sum_map_range_mono}
\label{List.sum_map_range_mono}
\leanok
\difficulty{4}
Let $g:\bbn\to\bbn$ take natural-number values, and let $a,b$ be natural numbers with
$a\le b$. The sum of $g(i)$ for $i<a$ is at most the sum of $g(i)$ for $i<b$.
\end{theorem}

\begin{theorem}[Successive partial sums of a list]
\lean{List.sum_take_succ_getD}
\label{List.sum_take_succ_getD}
\leanok
\difficulty{4}
Let $M$ be an additive monoid, let $l$ be a list of elements of $M$, and let $j$ be a
natural number. The sum of the first $j+1$ terms of $l$ equals the sum of its first $j$
terms plus the term at position $j$, taking that term to be $0$ when the list has no
such position.
\end{theorem}

\begin{theorem}[Successive partial sums of a mapped range]
\lean{List.sum_take_map_range_succ}
\label{List.sum_take_map_range_succ}
\leanok
\uses{List.getD_map_range, List.sum_take_succ_getD}
\difficulty{3}
Let $M$ be an additive monoid, let $f:\bbn\to M$, and let $n,j$ be natural numbers with
$j<n$. The sum of the first $j+1$ terms of the list $[f(0),f(1),\ldots,f(n-1)]$ equals
the sum of its first $j$ terms plus $f(j)$.
\end{theorem}

\begin{theorem}[Sum over a flattened list]
\lean{List.sum_flatMap_map}
\label{List.sum_flatMap_map}
\leanok
\difficulty{4}
Let $l$ be a finite list, let $g$ assign a finite list to each element of its type, and
let $F$ map the elements of those lists to integers. The sum of $F$ over the
concatenation of the lists $g(x)$ for $x\in l$ equals the sum, over $x\in l$, of the sum
of $F$ over $g(x)$.
\end{theorem}

\begin{theorem}[Sum over a filtered finite range]
\lean{List.sum_filter_finRange}
\label{List.sum_filter_finRange}
\leanok
\difficulty{3}
Let $m$ be a natural number, let $p$ be a predicate on the integers from $0$ through
$m-1$, and let $F$ assign an integer to each such integer. The sum of $F$ over the list
of elements of this range satisfying $p$ equals the sum of $F(x)$ over the finite set of
all such $x$.
\end{theorem}

\begin{theorem}[Absolute value of a sum]
\lean{List.abs_sum_le_sum_abs}
\label{List.abs_sum_le_sum_abs}
\leanok
\difficulty{1}
Let $R$ be a linearly ordered additive commutative group whose order is preserved by
addition. For every finite list $l$ of elements of $R$, the absolute value of its sum is
at most the sum of the absolute values of its entries.
\end{theorem}

\begin{theorem}[Absolute sum of an initial segment]
\lean{List.abs_sum_take_le_sum_abs}
\label{List.abs_sum_take_le_sum_abs}
\leanok
\uses{List.abs_sum_le_sum_abs}
\difficulty{4}
Let $R$ be a linearly ordered additive commutative group whose order is preserved by
addition, let $l$ be a finite list in $R$, and let $k\in\mathbb{N}$. The absolute value
of the sum of the first $k$ entries of $l$ (or all its entries if it has fewer than $k$)
is at most the sum of the absolute values of all entries of $l$.
\end{theorem}

\begin{theorem}[Bound for a truncated mapped sum]
\lean{List.abs_sum_take_map_le}
\label{List.abs_sum_take_map_le}
\leanok
\uses{List.abs_sum_take_le_sum_abs}
\difficulty{3}
Let $l$ be a finite list and let $f$ assign an integer to each element of its type.
Suppose there is an integer $B$ such that $|f(x)|\le B$ for every $x$ occurring in $l$.
For every natural number $k$, the absolute value of the sum of the first $k$ values of
$f$ on $l$ is at most $\operatorname{length}(l)B$.
\end{theorem}

\begin{theorem}[Bound on an integer sum]
\lean{List.abs_sum_map_le}
\label{List.abs_sum_map_le}
\leanok
\uses{List.abs_sum_take_map_le}
\difficulty{3}
Let $\alpha$ be a type, let $l$ be a finite list in $\alpha$, let
$f:\alpha\to\mathbb{Z}$, and let $B\in\mathbb{Z}$. If $|f(x)|\le B$ for every entry $x$
of $l$, then the absolute value of the sum of the values $f(x)$ over the list is at most
$\operatorname{length}(l)B$.
\end{theorem}

\begin{theorem}[List as the image of a finite set]
\lean{List.toFinset_eq_image_of_length_eq_card}
\label{List.toFinset_eq_image_of_length_eq_card}
\leanok
\difficulty{4}
Let $l$ be a duplicate-free list in a type with decidable equality, let $S$ be a finite
set, and let $g$ map the elements of $S$'s ambient type to the list's ambient type.
Suppose $g$ is injective, every entry of $l$ is the image under $g$ of some element of
$S$, and the length of $l$ equals the cardinality of $S$. Then the set of entries of $l$
is exactly the image of $S$ under $g$.
\end{theorem}

\begin{theorem}[Sum over an injective image of a finite set]
\lean{List.sum_map_eq_sum_of_length_eq_card}
\label{List.sum_map_eq_sum_of_length_eq_card}
\leanok
\uses{List.toFinset_eq_image_of_length_eq_card}
\difficulty{3}
Let $M$ be a commutative additive monoid, let $l$ be a list of elements of type $\beta$
with no repetitions, and let $S$ be a finite set of elements of type $\alpha$. Suppose
$g:\alpha\to\beta$ is injective, every element of $l$ is $g(x)$ for some $x\in S$, and
the length of $l$ equals the cardinality of $S$. Then the sum of $F$ over $l$ equals the
sum of $F(g(x))$ over $x\in S$, for any function $F:\beta\to M$.
\end{theorem}

\begin{theorem}[Bound for a finite running minimum]
\lean{List.foldl_min_le}
\label{List.foldl_min_le}
\leanok
\difficulty{4}
Let $\beta$ be a linearly ordered type, let $f:\mathbb{N}\to\beta$, and let $m$ be a
natural number. For every natural number $k\le m$, the minimum obtained by starting with
$f(0)$ and successively taking the minimum with $f(1),\ldots,f(m)$ is at most $f(k)$.
\end{theorem}

\begin{theorem}[A minimum is attained among finitely many values]
\lean{List.exists_foldl_min_eq}
\label{List.exists_foldl_min_eq}
\leanok
\difficulty{4}
Let $\beta$ be a linearly ordered type, let $f:\mathbb{N}\to\beta$, and let $m$ be a
natural number. There is an index $k\le m$ such that the minimum obtained by starting
with $f(0)$ and successively taking the minimum with $f(1),\ldots,f(m)$ equals $f(k)$.
\end{theorem}

\begin{theorem}[Count in a longer prefix]
\lean{List.count_take_succ}
\label{List.count_take_succ}
\leanok
\difficulty{3}
Let $l$ be a finite list, let $x$ be an element of its type, and let $i$ be a natural
number with $i<\operatorname{length}(l)$. The number of occurrences of $x$ among the
first $i+1$ entries equals its number among the first $i$ entries, plus $1$ if the entry
at position $i$ is $x$, and plus $0$ otherwise.
\end{theorem}

\begin{theorem}[Count in a longer prefix using a default]
\lean{List.count_take_succ_getD}
\label{List.count_take_succ_getD}
\leanok
\uses{List.count_take_succ}
\difficulty{3}
Let $l$ be a finite list, let $x$ and $d$ be elements of its type, and let $i$ be a
natural number with $i<\operatorname{length}(l)$. The number of occurrences of $x$ among
the first $i+1$ entries equals its number among the first $i$ entries, plus $1$ if the
entry at position $i$, using $d$ as the default when unavailable, is $x$, and plus $0$
otherwise.
\end{theorem}

\begin{theorem}[Prefix count is strictly smaller]
\lean{List.count_take_lt}
\label{List.count_take_lt}
\leanok
\uses{List.count_take_succ_getD}
\difficulty{4}
Let $l$ be a finite list, let $i$ be a natural number with $i<\operatorname{length}(l)$,
and let $x$ and $d$ be elements of its type. If the element at position $i$, using $d$
as the default when that position is unavailable, is $x$, then the number of occurrences
of $x$ among the first $i$ entries is strictly less than its number of occurrences in
$l$.
\end{theorem}

\begin{theorem}[Partitioning a threshold filter]
\lean{List.length_filter_lt_succ}
\label{List.length_filter_lt_succ}
\leanok
\difficulty{4}
Let $t$ assign a natural number to each element of a type, let $z$ be a natural number,
and let $l$ be a finite list of such elements. The number of entries of $l$ for which
$t(a)<z+1$ equals the number for which $t(a)<z$ plus the number for which $t(a)=z$.
\end{theorem}

\begin{theorem}[Counting a predicate on an initial interval]
\lean{List.length_filter_range}
\label{List.length_filter_range}
\leanok
\difficulty{1}
Let $q$ be a decidable predicate on the natural numbers and let $n$ be a natural number.
The number of indices $j<n$ satisfying $q(j)$ equals the count of numbers below $n$
satisfying $q$.
\end{theorem}

\begin{theorem}[Lookup by preceding satisfying indices]
\lean{List.getD_filter_range_count}
\label{List.getD_filter_range_count}
\leanok
\uses{List.length_filter_range}
\difficulty{5}
Let $q$ be a decidable predicate on the natural numbers, and let $n$ and $i$ be natural
numbers with $i<n$ and $q(i)$. In the list of numbers less than $n$ that satisfy $q$,
the entry at position equal to the number of indices less than $i$ satisfying $q$ is
$i$.
\end{theorem}

\begin{theorem}[Counting values in an indexed prefix]
\lean{List.count_take_ofFn}
\label{List.count_take_ofFn}
\leanok
\difficulty{4}
Let $n$ and $i$ be natural numbers, let $f$ assign an element of a type to each index
$j<n$, and let $a$ be an element of that type. The number of occurrences of $a$ among
the first $i$ entries of the list of values of $f$ equals the number of indices $j<n$
such that $j<i$ and $f(j)=a$.
\end{theorem}

\begin{theorem}[Counting values in an indexed list]
\lean{List.count_ofFn}
\label{List.count_ofFn}
\leanok
\uses{List.count_take_ofFn}
\difficulty{1}
Let $n$ be a natural number, let $f$ assign an element of a type to each index
$i\in\{0,\ldots,n-1\}$, and let $a$ be an element of that type. The number of
occurrences of $a$ in the list of values of $f$ equals the number of indices $i<n$ for
which $f(i)=a$.
\end{theorem}

\begin{theorem}[Membership after dropping initial smaller elements]
\lean{List.mem_dropWhile_lt}
\label{List.mem_dropWhile_lt}
\leanok
\difficulty{4}
Let $\alpha$ be a linearly ordered type, let $a,c\in\alpha$, and let $l$ be a list whose
entries are pairwise strictly increasing. Then $c$ belongs to the list obtained by
dropping the initial entries of $l$ that are less than $a$ if and only if $c$ belongs to
$l$ and $a\le c$.
\end{theorem}

\begin{theorem}[Membership among initial smaller elements]
\lean{List.mem_takeWhile_lt}
\label{List.mem_takeWhile_lt}
\leanok
\difficulty{4}
Let $\alpha$ be a linearly ordered type, let $b,c\in\alpha$, and let $l$ be a list whose
entries are pairwise strictly increasing. Then $c$ belongs to the initial segment of $l$
consisting of entries less than $b$ if and only if $c$ belongs to $l$ and $c<b$.
\end{theorem}

\begin{theorem}[A new value outside the preceding prefix]
\lean{List.getD_notMem_take_of_sorted}
\label{List.getD_notMem_take_of_sorted}
\leanok
\difficulty{5}
Let $l$ be a finite list in a linearly ordered set, with each entry at most the next.
For an index $i<\operatorname{length}(l)$ and a default value $d$, suppose either $i=0$
or the entry at index $i$ differs from the entry at index $i-1$, using $d$ for any
out-of-range index. Then the entry at index $i$, again using $d$ as the default, does
not occur among the first $i$ entries of $l$.
\end{theorem}

\begin{definition}[Uniform absolute bound]
\lean{Light.AbsLe}
\label{Light.AbsLe}
\leanok
A list $l$ of integers is bounded in absolute value by an integer $U$ if every entry $x$
of $l$ satisfies $|x|\le U$.
\end{definition}

\begin{theorem}[Bound on an in-range list entry]
\lean{Light.AbsLe.getElem}
\label{Light.AbsLe.getElem}
\leanok
\uses{Light.AbsLe}
\difficulty{1}
Let $l$ be a list of integers and let $U$ be an integer. If every entry of $l$ has
absolute value at most $U$, then for every natural number $i$ with
$i<\operatorname{length}(l)$, the entry at position $i$ has absolute value at most $U$.
\end{theorem}

\begin{theorem}[Absolute bounds are preserved under concatenation]
\lean{Light.AbsLe.append}
\label{Light.AbsLe.append}
\leanok
\uses{Light.AbsLe}
\difficulty{1}
Let $l_1$ and $l_2$ be lists of integers and let $U$ be an integer. If every entry of
each list has absolute value at most $U$, then every entry of their concatenation also
has absolute value at most $U$.
\end{theorem}

\begin{theorem}[Bound on a defaulted list entry]
\lean{Light.AbsLe.abs_getD_le}
\label{Light.AbsLe.abs_getD_le}
\leanok
\uses{Light.AbsLe, List.getD_of_forall_mem}
\difficulty{1}
Let $l$ be a list of integers and $U$ an integer with $0\le U$. If every entry $x$ of
$l$ satisfies $|x|\le U$, then for every natural number $i$, the value at position $i$
with default value zero also satisfies $|l_i^{(0)}|\le U$.
\end{theorem}

\begin{definition}[Coordinatewise sum of lists]
\lean{Light.sumLists}
\label{Light.sumLists}
\leanok
Given a natural number $m$ and a list of lists of integers, start with the list of $m$
zeros and successively replace the current list by the pairwise sum with each list,
truncating each pairwise sum to the length of the shorter list. The resulting list is
the defined sum.
\end{definition}

\begin{theorem}[Coordinatewise sum of mapped lists]
\lean{Light.sumLists_map}
\label{Light.sumLists_map}
\leanok
\uses{Light.sumLists}
\difficulty{6}
Let $l$ and $ts$ be finite lists, and let $f$ assign an integer to each pair consisting
of an element of $ts$ and an element of $l$. Adding the lists obtained by mapping
$f(t,\cdot)$ over $l$, coordinate by coordinate and starting with a zero list of length
$|l|$, gives the list whose entry at $c\in l$ is $\sum_{t\in ts} f(t,c)$.
\end{theorem}
